% Data Archiving System TRS; WBS 123456.02.01.
\section{Process for Establishing a Technical Requirements Specification}

\subsection{General}

The establishment of the Data Archiving System (DAS) Technical Requirements Specification (TRS) follows a progressive requirements definition process. The process begins with the identification of the mission and user needs and gradually develops into a structured and verifiable set of technical requirements.

At the initial stage, the requirements are defined at a high level and describe the intended functions, capabilities and constraints of the DAS. As the project progresses and the system concept is further developed, these requirements are reviewed, refined and supplemented with more detailed requirements concerning data acquisition, storage, organization, preservation, retrieval, security and integrity.

Within the EWAS project, the DAS TRS represents the technical requirements applicable to both the Raw Data Archive (RDA) and the Processed Data Archive (PDA). It provides the basis for the subsequent design and development of the archiving functions within the Ground Segment.

The requirements are established from the perspective of the customer and are intended to describe the problem to be solved and the expected system capabilities, rather than prescribe a particular technical implementation. The selection of specific technologies, storage solutions or software architectures is therefore left to the subsequent system design and development activities.

The DAS TRS is intended to serve as a technical reference for the development, verification and acceptance of the Data Archiving System. Any changes to the agreed requirements shall be handled through the applicable project requirements and configuration management processes.

\subsection{Process for Establishing the DAS Technical Requirements Specification}

The process for establishing the DAS TRS is based on the progressive identification, structuring, assessment and refinement of technical requirements. It is adapted to the development stage and scope of the EWAS project and is applicable to the DAS as a Ground Segment subsystem.

The process starts with the identification of the needs arising from the EWAS mission objectives, Ground Segment operations and the expected users of the archived data. These needs are translated into technical requirements describing the capabilities expected from the DAS.

The initial requirements are then reviewed and structured to ensure that they provide a complete and consistent description of the system's intended functions. As the DAS concept develops, the requirements are reassessed to account for the selected system concept, identified interfaces and applicable operational constraints.

The process consists of the following activities:

\begin{description}
    \item[\textbf{F1.1 -- Identification and capture of needs:}] The customer identifies the needs and expectations of the DAS stakeholders, including Ground Segment operators, authorised crisis-management authorities, analysts and technical and maintenance personnel. The activity considers the mission objectives, operational use of archived data, data sources, expected data products and applicable environmental, security and programme constraints.
    \item[\textbf{F1.2 -- Structuring, classification and justification of requirements:}] The identified needs are translated into individual technical requirements. Requirements are organised according to the functional areas of the DAS, including data ingestion, raw data archiving, processed data archiving, data retrieval, metadata management, data retention, data integrity, access control, backup and recovery. Each requirement is justified by its relationship to an identified stakeholder need or system function.
    \item[\textbf{F1.3 -- Assessment of the requirements set:}] The complete set of requirements is assessed to ensure correctness, consistency, completeness and suitability for the intended use. Particular attention is given to the separation of RDA and PDA responsibilities, the traceability of archived data and the distinction between archival functions and data processing or operational decision-making functions. Requirements are also assessed for clarity, achievability and verifiability.
    \item[\textbf{F1.4 -- Establishment and release of the preliminary DAS TRS:}] The assessed requirements are consolidated into a preliminary version of the DAS TRS. This version documents the initial system scope, expected capabilities, operational constraints and technical requirements. It provides the baseline for further development of the DAS concept.
    \item[\textbf{F1.5 -- Review of the preliminary TRS and identification of possible concepts:}] The preliminary DAS TRS is reviewed against possible system concepts. Alternative approaches to implementing the RDA and PDA functions are considered, including their implications for data organisation, storage capacity, data retrieval, security, maintainability and interfaces with other Ground Segment elements.
    \item[\textbf{F1.6 -- Evaluation and selection of the preferred concept:}] The identified concepts are evaluated against the mission needs and the requirements defined in the preliminary TRS. The preferred concept is selected based on its ability to satisfy the requirements and its compatibility with the overall EWAS Ground Segment architecture.
    \item[\textbf{F1.7 -- Review and adjustment of technical requirements:}] The preliminary requirements are reviewed in light of the selected concept. Requirements are adjusted, clarified or added where necessary to reflect the capabilities and constraints identified during concept evaluation. This activity includes the refinement of requirements related to data retention, historical data preservation, data traceability and the support of future analysis and algorithm evaluation activities.
    \item[\textbf{F1.8 -- Final structuring, classification and justification of requirements:}] The updated requirements are organised into a consistent structure. Each requirement is assigned to the appropriate DAS functional area and checked for clear identification, justification and traceability to the relevant stakeholder need or higher-level system requirement.
    \item[\textbf{F1.9 -- Final assessment of the requirements set:}] The complete DAS TRS is assessed to confirm that the requirements are correct, consistent, sufficiently complete and suitable for the intended system scope. The assessment also verifies that requirements are individually understandable and can be verified through an appropriate verification method.
    \item[\textbf{F1.10 -- Establishment and release of the DAS TRS:}] Following the assessment, the DAS TRS is consolidated and formally released. The released specification provides the technical requirements baseline for the subsequent design, implementation and verification activities of the Data Archiving System.
\end{description}

The process is iterative. Requirements may be reviewed and updated as a result of design development, interface definition, verification planning or changes to the overall EWAS system requirements. All changes to the approved DAS TRS shall be subject to the applicable change control and configuration management procedures.

The process described above applies to the DAS at its subsystem level and may also be applied to its constituent elements, namely the Raw Data Archive and the Processed Data Archive, where more detailed requirements specifications are required.

The outcome of the process is the DAS TRS, which defines the technical requirements to be satisfied by the Data Archiving System. The specification provides the basis for the development agreement and for subsequent design, verification and acceptance activities.
